\clearpage
\VerdierPhysicalPageAnchor{0236}%
\begingroup\small\setlength{\parskip}{0pt}\setlength{\emergencystretch}{1.5em}
\setlength{\abovedisplayskip}{4pt plus 2pt minus 2pt}%
\setlength{\belowdisplayskip}{4pt plus 2pt minus 2pt}%

\VerdierSemanticAnchor{ch4d-4-2-4}
\noindent\textbf{4.2.4.}\quad
A morphism between filtered objects of type $\widetilde{\mathbb Z}$ in
$\operatorname{comp}(\mathcal A)$ induces, in the evident way, a morphism
between the associated spectral objects described in Proposition (4.2.3),
a). Proposition (4.2.3) therefore defines a functor from the category of
filtered objects of $\operatorname{comp}(\mathcal A)$ to the category of
spectral objects with values in $D(\mathcal A)$.

\VerdierSemanticAnchor{ch4d-4-2-4-first-canonical}
Let $X$ be an object of $\operatorname{comp}(\mathcal A)$ and let
$X(q)\subset X$, $q\in\mathbb Z$, be its first canonical filtration. In
this case, the spectral object associated by Proposition (4.2.3) with the
filtered object $\bigl(X,X(q),q\in\mathbb Z\bigr)$ is precisely the first
spectral object, with values in $D(\mathcal A)$, associated with the
complex $X$ (4.1.5).

\VerdierSemanticAnchor{ch4d-4-2-4-nfold}
Likewise, let $X$ be an $n$-fold complex in $\mathcal A$, let
$\int_n X=Y$ be the associated simple complex, let $m$ be an integer,
$1\leq m\leq n$, and let
\[
0={}^{m}\!Y(-\infty)\subset\cdots\subset{}^{m}\!Y(q)
\subset\cdots\subset{}^{m}\!Y(+\infty)=Y
\]
be the filtration defined by the $m$-th degree of $X$ (4.1.6.9). In this
case, the spectral object associated by Proposition (4.2.3) with the
filtered object $\bigl(Y,{}^{m}\!Y(q),q\in\mathbb Z\bigr)$ is precisely
the spectral object, with values in $D(\mathcal A)$, defined by $X$
filtered by its $m$-th degree (4.1.8).

\VerdierSemanticAnchor{ch4d-4-3}
\medskip\noindent\textbf{4.3. Second canonical spectral object.}

\VerdierSemanticAnchor{ch4d-4-3-1}
\medskip\noindent\textbf{4.3.1.}\quad
Let $X$ be an object of $\operatorname{comp}(\mathcal A)$. Set

\VerdierSemanticAnchor{ch4d-4-3-1-1}
\begin{equation}
\begin{aligned}
X(-\infty)&=0,\qquad X(+\infty)=X,\\
X(q)&=\cdots\longrightarrow X^{q-3}\longrightarrow X^{q-2}
\longrightarrow\operatorname{Ker}(d_X^{q-1})
\longrightarrow0\longrightarrow0\longrightarrow\cdots,
\qquad q\in\mathbb Z,
\end{aligned}
\tag{4.3.1.1}
\end{equation}
the differentials of the latter complex being induced by the differentials
of $X$. There are canonical monomorphisms

\VerdierSemanticAnchor{ch4d-4-3-1-2}
\begin{equation}
0=X(-\infty)\subset\cdots\subset X(q)\subset X(q+1)
\subset\cdots\subset X(+\infty)=X
\tag{4.3.1.2}
\end{equation}
that define a filtration of type $\widetilde{\mathbb Z}$ on $X$, called
the \emph{second canonical filtration of $X$}. Proposition (4.2.3) then
yields a spectral object $\bigl(X(p,q),\delta(p,q,r)\bigr)$, called the
\emph{second spectral object, with values in $D(\mathcal A)$, associated
with the complex $X$}.

\VerdierSemanticAnchor{ch4d-4-3-1-associated-functor}
A morphism in $\operatorname{comp}(\mathcal A)$,
$X\longrightarrow Y$, respects the second canonical
filtrations; consequently, the functor
\[
X\longmapsto\bigl(X(p,q),\delta(p,q,r)\bigr)
\]
is a functor from $\operatorname{comp}(\mathcal A)$ to the category of
spectral objects with values in $D(\mathcal A)$. This functor is called the
\emph{second canonical spectral-object functor}.
\endgroup
